\subsection{HBN Isomers}

Having validated the tensor reconstruction procedure for a molecule
with strongly representation-dependent Watson constants, we now apply
the tensor analysis to a series of related molecules in order to
investigate the structural trends encoded in the reduced quartic
tensor itself.

The rotational constants and quartic centrifugal distortion constants
of the four HBN isomers (\emph{o}-HBN, \emph{p}-HBN, \emph{cis}-m-HBN,
and \emph{trans}-m-HBN) computed at the HPCS2 and DPCS3 levels are
reported in Table~\ref{table:new_results} together with the
experimental parameters obtained from the spectroscopic analysis.

In addition to the Watson centrifugal distortion constants, the table
reports the diagnostic parameter $s_{111}$ computed in the $I^r$
representation (A reduction). For compactness the values are given as
$\log_{10}|s_{111}|$. The table further includes constants obtained
through representation transformations from $I^r$ (A) to $III^r$ (A)
and to $I^r$ (S) using the tensor transformation procedure described
above.

The values of $s_{111}$ provide useful information on the numerical
conditioning of the effective rotational Hamiltonian under different
reductions and axis representations. Several systematic trends emerge
from Table~\ref{table:new_results}.

Within the present implementation these gauge diagnostics are most
useful when considered together with the condition numbers of the
representation transforms.  In this way a single parameter set may be
transported to alternative reductions and axis representations,
examined in terms of its tensor invariants, gauge sensitivity, and
numerical conditioning, and only then selected as the starting point
for any new least-squares refinement.  The tensor transform therefore
acts not only as a formal change of representation but also as a
practical prescreening tool for spectroscopic analysis.

\begin{table*}[ht]
  \centering
  \scriptsize
  \setlength{\tabcolsep}{5pt}

  \caption{Ground-state rotational constants (MHz) and quartic
  centrifugal distortion constants (kHz) for the four HBN isomers.
  Ground-state constants were obtained by subtracting the HPCS2
  vibrational corrections from the equilibrium values.
  The HPCS2 vibrational corrections are
  $\Delta B_a^{vib}=23.6$, $6.2$, and $5.5$ MHz for
  cis-$o$-HBN,
  $42.3$, $4.2$, and $4.0$ MHz for $p$-HBN,
  $20.2$, $5.6$, and $4.5$ MHz for cis-$m$-HBN,
  and $22.4$, $5.1$, and $4.4$ MHz for trans-$m$-HBN
  (for $B_a$, $B_b$, and $B_c$, respectively).
  Additional rows report the gauge diagnostic parameter $s_{111}$
  computed in $I^r$ (A reduction), together with constants transformed
  from $I^r$ (A) to $III^r$ (A) and to $I^r$ (S) through the tensor
  pipeline described in this work. For compactness,
  $s_{111}$ values are reported as $\log_{10}|s_{111}|$.}

  \begin{tabular}{lcccccccccccc}
  \toprule

  & \multicolumn{3}{c}{\textbf{cis-$o$-HBN}}
  & \multicolumn{3}{c}{\textbf{$p$-HBN}}
  & \multicolumn{3}{c}{\textbf{cis-$m$-HBN}}
  & \multicolumn{3}{c}{\textbf{trans-$m$-HBN}} \\

  Parameter
  & Exp. & HPCS2 & DPCS3
  & Exp. & HPCS2 & DPCS3
  & Exp. & HPCS2 & DPCS3
  & Exp. & HPCS2 & DPCS3 \\

  \midrule
  \multicolumn{13}{c}{\textbf{Ground-state rotational constants (MHz)}} \\

  $B_a$ &
  3053.7 & 3037.2 & 3062.1 &
  5613.0 & 5596.5 & 5638.8 &
  3408.9 & 3398.4 & 3418.6 &
  3403.1 & 3390.4 & 3412.8 \\

  $B_b$ &
  1511.3 & 1506.4 & 1512.6 &
  990.5 & 986.5 & 990.7 &
  1205.8 & 1201.0 & 1206.6 &
  1208.5 & 1204.1 & 1209.2 \\

  $B_c$ &
  1010.8 & 1006.8 & 1012.3 &
  841.9 & 838.7 & 842.7 &
  890.7 & 887.3 & 891.8 &
  891.8 & 888.4 & 892.9 \\

  \midrule
  \multicolumn{13}{c}{\textbf{Quartic centrifugal distortion constants in $I^r$ (A reduction) (kHz)}} \\

  $\Delta_J$ &
  0.038 & 0.0368 & 0.0370 &
  0.020 & 0.015 & 0.016 &
  0.041 & 0.0387 & 0.0390 &
  0.040 & 0.0380 & 0.0383 \\

  $\Delta_{JK}$ &
  0.624 & 0.619 & 0.6204 &
  0.231 & 0.221 & 0.223 &
  0.027 & 0.0246 & 0.0250 &
  0.020 & 0.0352 & 0.0356 \\

  $\Delta_K$ &
  -0.100 & 0.381 & 0.384 &
  9 & 0.918 & 0.922 &
  1.180 & 1.173 & 1.1778 &
  1.130 & 1.158 & 1.161 \\

  $\delta_J$ &
  0.011 & 0.010 & 0.010 &
  0.005 & 0.003 & 0.003 &
  0.014 & 0.014 & 0.014 &
  0.014 & 0.014 & 0.014 \\

  $\delta_K$ &
  0.280 & 0.332 & 0.335 &
  0.400 & 0.182 & 0.184 &
  0.160 & 0.182 & 0.185 &
  0.220 & 0.182 & 0.185 \\

  $\log_{10}|s_{111}|$ &
  -10.43 & -10.33 & -10.33 &
  -11.27 & -11.71 & -11.71 &
  -11.36 & -11.20 & -11.19 &
  -11.00 & -11.19 & -11.18 \\

  \midrule
  \multicolumn{13}{c}{\textbf{Transformed constants in $III^r$ (A reduction) (kHz)}} \\

  $\Delta_J(III^r)$ &
  0.5310 & 0.6411 & 0.6445 &
  1.9072 & 0.4028 & 0.4068 &
  0.3718 & 0.3818 & 0.3851 &
  0.3990 & 0.3839 & 0.3869 \\

  $\Delta_{JK}(III^r)$ &
  -0.9650 & -1.2412 & -1.2499 &
  -4.0638 & -0.8502 & -0.8572 &
  -0.8222 & -0.8700 & -0.8788 &
  -0.9420 & -0.8703 & -0.8785 \\

  $\Delta_K(III^r)$ &
  0.4500 & 0.6168 & 0.6223 &
  2.1667 & 0.4563 & 0.4603 &
  0.4633 & 0.4988 & 0.5046 &
  0.5550 & 0.4963 & 0.5018 \\

  $\delta_J(III^r)$ &
  0.1255 & 0.2450 & 0.2461 &
  2.3052 & 0.2832 & 0.2847 &
  0.2947 & 0.2924 & 0.2937 &
  0.2805 & 0.2913 & 0.2921 \\

  $\delta_K(III^r)$ &
  0.1150 & 0.2613 & 0.2635 &
  2.4500 & 0.3205 & 0.3225 &
  0.3750 & 0.3843 & 0.3870 &
  0.3925 & 0.3805 & 0.3828 \\

  $\log_{10}|s_{111}|$ &
  -10.12 & -9.77 & -9.78 &
  -8.75 & -9.63 & -9.63 &
  -9.58 & -9.57 & -9.57 &
  -9.57 & -9.57 & -9.58 \\

  \midrule
  \multicolumn{13}{c}{\textbf{Transformed constants in $I^r$ (S reduction) (kHz)}} \\

  $D_J$ &
  0.0380 & 0.0368 & 0.0370 &
  0.0200 & 0.0150 & 0.0160 &
  0.0410 & 0.0387 & 0.0390 &
  0.0400 & 0.0380 & 0.0383 \\

  $D_{JK}$ &
  0.6240 & 0.6190 & 0.6204 &
  0.2310 & 0.2210 & 0.2230 &
  0.0270 & 0.0246 & 0.0250 &
  0.0200 & 0.0352 & 0.0356 \\

  $D_K$ &
  -0.1000 & 0.3810 & 0.3840 &
  9.0000 & 0.9180 & 0.9220 &
  1.1800 & 1.1730 & 1.1778 &
  1.1300 & 1.1580 & 1.1610 \\

  $d_1$ &
  0.5820 & 0.6840 & 0.6900 &
  0.8100 & 0.3700 & 0.3740 &
  0.3480 & 0.3920 & 0.3980 &
  0.4680 & 0.3920 & 0.3980 \\

  $d_2$ &
  -0.5600 & -0.6640 & -0.6700 &
  -0.8000 & -0.3640 & -0.3680 &
  -0.3200 & -0.3640 & -0.3700 &
  -0.4400 & -0.3640 & -0.3700 \\

  $\log_{10}|s_{111}|$ &
  -9.33 & -9.26 & -9.26 &
  -9.22 & -9.55 & -9.55 &
  -9.55 & -9.50 & -9.50 &
  -9.42 & -9.50 & -9.50 \\

  \bottomrule
  \end{tabular}

  \label{table:new_results}

\end{table*}

Within the present tensor framework, the quartic Watson constants are
best regarded as coordinates obtained by projection of the reduced
quartic tensor $\Tau'$.  The physically meaningful,
representation-independent descriptors are therefore the invariants of
$\Tau'$ itself rather than the individual Watson constants.

For a symmetric $3\times 3$ tensor, the most natural invariants are
either the three eigenvalues
\[
\lambda_1,\lambda_2,\lambda_3,
\]
or, equivalently, the three polynomial invariants
\begin{align}
I_1^{(\Tau')} &= \mathrm{Tr}(\Tau'),\\
I_2^{(\Tau')} &= \tfrac12\!\left[
\bigl(\mathrm{Tr}(\Tau')\bigr)^2-\mathrm{Tr}\bigl((\Tau')^2\bigr)
\right],\\
I_3^{(\Tau')} &= \det(\Tau').
\end{align}
These quantities remain unchanged under permutation similarity,
\[
\Tau' \longrightarrow P^T\Tau' P,
\]
and therefore provide the appropriate basis for comparing different
representations and different molecules.

To illustrate the physical content of this analysis we applied the
tensor reconstruction procedure to the four HBN isomers using the HPCS2
quartic constants.  The resulting spectral invariants of $\Tau'$ are
reported in Table~\ref{table:hbn_tauprime}.  In contrast to the Watson
constants, which may vary substantially with the chosen reduction or
axis representation, these quantities characterize the underlying
quartic centrifugal distortion field in a representation-independent
way.

\begin{table}[h]
\centering
\caption{Spectral invariants of the reduced quartic tensor $\Tau'$ for
the HBN isomers.  The numerical values should be obtained from the
reconstructed tensor in the chosen reference representation.}

\begin{tabular}{lcccccc}
\hline
Isomer & $\lambda_1$ & $\lambda_2$ & $\lambda_3$ & $I_1^{(\Tau')}$ & $I_2^{(\Tau')}$ & $I_3^{(\Tau')}$ \\
\hline
cis-$o$-HBN   & 0.9585 & -0.2946 & -5.1055 & -4.4416 & -3.6721 & 1.4417 \\
$p$-HBN       & 0.0267 & -0.0328 & -4.7299 & -4.7360 & 0.0280 & 0.0042 \\
cis-$m$-HBN   & -0.0243 & -0.2578 & -4.9727 & -5.2548 & 1.4090 & -0.0311 \\
trans-$m$-HBN & -0.0214 & -0.2508 & -4.9566 & -5.2288 & 1.3546 & -0.0266 \\
\hline
\end{tabular}

\label{table:hbn_tauprime}

\end{table}

A clear structural pattern is expected to emerge from these tensorial
invariants.  In particular, the two meta isomers should exhibit very
similar spectral invariants if their intrinsic quartic centrifugal
distortion fields are nearly the same, whereas the ortho and para
isomers may display a different balance of principal tensor components.
Such comparisons are far more robust at the level of $\Tau'$ than at
the level of the Watson constants themselves, because the latter are
strongly affected by the chosen parametrization of the effective
Hamiltonian.

From this point of view, the HBN series illustrates the principal
advantage of the tensor formulation.  Watson constants are convenient
for direct spectral fitting, but they mix intrinsic tensorial content
with representation-dependent parametrization.  The invariants of
$\Tau'$, by contrast, provide a compact and representation-independent
description of the underlying quartic centrifugal distortion field and
therefore furnish a more suitable basis for comparing related
molecules, computational models, and experimental analyses.  At the
same time, the associated values of $s_{111}$ and the transformation
condition numbers provide a practical criterion for deciding which
representation is numerically most appropriate before repeating an
explicit spectroscopic fit.
